\chapter{Background and Related Work}\label{cpt:background}

This chapter positions the dissertation within four bodies of work: code-generation models and executable benchmarks, watermarking for generated text and code, data provenance and contamination analysis, and security evaluation for generated code. The purpose is not to survey every paper in these areas. Instead, the chapter identifies the assumptions that motivate WatermarkScope's design: code artifacts must execute, benchmark claims need stable denominators, watermark detection needs negative controls, and security-sensitive systems need abstention rather than forced decisions.

\section{Code Generation Models}

Large language models trained on code have made program synthesis a practical research area. Earlier neural code models such as CodeBERT, PLBART, and CodeT5 established reusable code representations and sequence-to-sequence pre-training for program understanding and generation~\cite{feng2020codebert,ahmad2021plbart,wang2021codet5}. Codex and HumanEval then established pass@k-style functional evaluation for generated Python programs~\cite{chen2021codex}. Later generative code models such as InCoder and Code Llama broadened the setting to infilling, editing, and local open-model evaluation~\cite{fried2023incoder,roziere2023codellama}. These developments motivate this dissertation's model setting: source-code watermarking should be evaluated on executable model outputs, not only on token streams.

The executable nature of code changes the meaning of watermark quality. A generated answer can carry a watermark while failing tests, or pass tests while losing the watermark after ordinary edits. A useful evaluation must therefore track both evidence quality and program utility. This principle appears throughout CodeMarkBench and the later protocol modules.

\section{Executable Code Benchmarks}

Code evaluation benchmarks differ in scope. CodeXGLUE aggregates a broad set of code understanding and generation tasks~\cite{lu2021codexglue}. MBPP focuses on short Python programming problems~\cite{austin2021mbpp}, while APPS introduces competition-style programming tasks with larger solution spaces~\cite{hendrycks2021apps}. MultiPL-E extends code-generation evaluation across programming languages~\cite{cassano2022multiple}. DS-1000 emphasizes data-science code and highlights the difficulty of contamination-aware evaluation~\cite{lai2023ds1000}. SWE-bench moves toward repository-level issue resolution~\cite{jimenez2024swebench}, and LiveCodeBench explicitly targets live benchmark construction to reduce contamination risks~\cite{jain2024livecodebench}. BigCodeBench further broadens execution-based code evaluation to library-rich tasks requiring multiple function calls and complex instructions~\cite{zhuo2025bigcodebench}.

These benchmarks motivate several design decisions in CodeMarkBench. First, one source group is not sufficient; a benchmark should include public and crafted sources. Second, language coverage matters because transformations that preserve Python behavior may not preserve Java, C++, Go, or JavaScript behavior. Third, a benchmark should record exact run identities and release contracts, because later comparisons are meaningful only if the denominator remains stable.

\section{Watermarking Generated Text and Code}

Watermarking for generated text often uses a keyed bias in token selection and a statistical detector over the generated sequence. Kirchenbauer et al. introduced a widely cited green-list watermarking formulation for large language models~\cite{kirchenbauer2023watermark}. Subsequent work studies robustness, reliability, provability, and undetectability under different edit or access assumptions~\cite{zhao2024provable,kirchenbauer2024reliability,christ2024undetectable}. Code watermarking, however, has additional structure. Generated code contains syntax trees, control flow, data flow, dependencies, tests, comments, and naming conventions. Code-oriented watermarking and API watermarking work therefore needs to balance detector strength with executable utility and edit robustness~\cite{li2023protecting,lee2024who,suresh2024robust}.

This dissertation uses prior watermarking work in two ways. CodeMarkBench evaluates existing runtime code-watermarking baselines under a shared release contract whose repository and artifact entry points are recorded in Appendix~\ref{app:repository-access}. SemCodebook then moves away from purely surface-level evidence by treating structured program carriers as recoverable provenance support. This does not make the watermark universal; it makes the claim more explicit about carrier families and admitted model cells.

\section{Why Code Changes the Watermarking Problem}

The difference between text and code watermarking is not only that code has syntax. The deeper issue is that ordinary software maintenance can change the observable surface while preserving the intended behavior. A developer may rename variables, extract a helper, remove comments, normalize formatting, replace an expression with an equivalent library call, or translate a solution into another language. A detector that depends on the original token surface can fail after these edits even when the program remains a legitimate descendant of the generated artifact.

The reverse failure is also possible. A detector can fire on clean code because a formatting style, naming pattern, or common control-flow idiom happens to resemble a watermark carrier. This is especially dangerous for source-code watermarking because false positives can imply authorship, contamination, or security risk. WatermarkScope's design therefore treats negative controls and executable preservation as part of the central evidence, not as a secondary sanity check. A result that has a high detector rate but an unclear control denominator would be weak evidence for this dissertation.

This observation explains why the later modules have different access assumptions. SemCodebook uses structured carriers because white-box access allows the method to inspect program representations rather than only token strings. CodeDye does not attempt the same recovery in a black-box API setting because the necessary carriers and model internals are unavailable. ProbeTrace adds an owner registry because detection is not the same as attribution. SealAudit adds a triage resolver because provenance evidence does not answer whether the watermarking mechanism is safe.

\section{Provenance, Memorization, and Contamination}

Training data extraction and memorization studies show that large language models can reproduce or expose training examples under certain conditions~\cite{carlini2021extracting,carlini2023quantifying}. Membership and pretraining-data detection methods, such as Min-K\% probability approaches, study whether a model has likely seen a candidate sequence~\cite{shi2023detecting}. Benchmark contamination surveys further show that evaluation datasets can leak into pretraining or instruction-tuning corpora, complicating claims about model capability~\cite{xu2024contamination}. Data statements and dataset documentation encourage explicit provenance, collection, and intended-use descriptions~\cite{bender2018data}.

These works motivate CodeDye's conservative framing. A black-box audit may surface sparse evidence, but it should not turn a small signal count into a contamination prevalence estimate. Instead, it should preserve prompt hashes, raw response hashes, structured payload hashes, controls, and threshold versions so that later human or provider-side review can examine the evidence without changing the denominator.

The same provenance literature also clarifies why attribution is not the same as detection. A detector can say that an output resembles a class of examples, while an attribution protocol must bind the output to a registered source or owner under a precommitted rule. ProbeTrace therefore introduces an active owner registry rather than treating high similarity as authorship. This distinction is important for the dissertation story: CodeDye preserves black-box evidence when ownership is unknown, while ProbeTrace studies the stronger but narrower setting where the owner key and source split have already been registered.

\section{Security Evaluation for Generated Code}

Generated-code safety research shows that code assistants can produce insecure code or influence developers toward insecure patterns. Pearce et al. studied insecure code suggestions from GitHub Copilot~\cite{pearce2022asleep}. Perry et al. found that AI assistants can affect whether users write secure or insecure code~\cite{perry2023users}. LLMSecEval provides a dataset for evaluating security properties of LLM-generated code~\cite{siddiq2023llmseceval}. Security taxonomies such as the MITRE CWE Top 25 provide practical categories for common software weaknesses~\cite{mitre2025cwe}.

These studies motivate SealAudit. A watermark or marker should not be treated only as a provenance signal. If a marker changes model behavior, triggers unsafe completions, or interacts with laundering and spoofability attacks, it becomes a security-relevant object. SealAudit therefore reports selective triage with explicit needs-review outcomes rather than claiming that a model is safe.

\section{Gap Addressed by This Dissertation}

The gap is not simply the absence of another detector. The gap is an evidence protocol for deciding what a detector result is allowed to mean. The works reviewed above motivate this point from different directions: code benchmarks emphasize executable evaluation, watermarking papers emphasize keyed detection and robustness, provenance work warns against unsupported data-use inference, and security studies show that generated code can carry deployment risk. WatermarkScope addresses the combined gap by putting benchmark foundation, structured provenance recovery, conservative null-audit, active attribution, and selective triage under one claim-boundary discipline.

A second gap is methodological rather than only empirical. Many watermarking papers optimize a detector score, but a source-code watermarking system also has to specify what counts as a valid observation. A syntactic perturbation that changes the test outcome should not be treated the same as an executable semantic-preserving rewrite. A negative control that lacks the target watermark must be preserved even when it makes the result less impressive. A row generated for debugging, canary checking, or failure localization should not silently enter a main table. This dissertation therefore treats evidence admission as part of the method. The admission rule defines which rows may support a claim before the detector result is known.

The final gap concerns how to compare results across different access assumptions. White-box provenance, black-box null-audit, active-owner attribution, and security triage are not interchangeable tasks. A high recovery rate in an admitted white-box setting does not imply contamination detection in a black-box provider setting. Conversely, a sparse black-box audit signal is not a failed watermark if the stated goal is to preserve auditable evidence under conservative thresholds. The WatermarkScope framing is useful because it makes these access assumptions explicit. Each module reports the strongest claim that its evidence supports, and the dissertation avoids converting one module's metric into another module's conclusion.

These gaps appear sequentially as access changes: benchmark, provenance, audit, attribution, and triage. The rest of the dissertation follows that sequence instead of treating the artifacts as independent endpoints. Code-generation benchmarks contribute executable tasks and contamination-aware evaluation pressure; watermarking contributes keyed evidence and robustness concerns; provenance work contributes documentation and auditability; security work contributes the idea that abstention can be safer than forced classification. WatermarkScope's novelty is to make these standards interact. A row that is executable but not hash-bound is insufficient for black-box audit. A detector hit that is not owner-bound is insufficient for attribution. A marker-hidden case with conflicting safety evidence is insufficient for automatic acceptance.

This synthesis also narrows the dissertation's novelty claim. WatermarkScope does not replace HumanEval-style evaluation, green-list watermarking, contamination detection, attribution studies, or generated-code safety benchmarks. Instead, it connects their evaluation lessons into a source-code watermarking lifecycle with fixed denominators, explicit controls, uncertainty, support-only ledgers, and access-specific claim boundaries. The design test for each stage is therefore operational: can an examiner identify the admitted denominator, the control family, the access model, the uncertainty statement, and the forbidden interpretation? If any of these elements is missing, a numerical result becomes hard to defend even when the numerator is large.

This is the main reason the dissertation begins with a benchmark rather than with the strongest watermarking result. Before a provenance method can claim recovery, the report must define which executable artifacts are eligible for recovery. Before a black-box audit can report sparse signals, it must define which transcripts and controls are bound to the audit surface. Before attribution or triage can be interpreted, the reader must know whether source ownership or safety risk was actually the question being tested. CodeMarkBench is introduced next because it supplies this first stable denominator and establishes the counting discipline used by the later modules.

